\section{Posterior-flow representation}
\label{sec:flow}

The separator must control a ratio of Gaussian mixtures without replacing
its shared numerator and denominator by independent intervals.  This section
identifies a small posterior state that retains the relevant dependence.

For an observation $y$, let
\[
 p_j(y)=\frac{w_j\exp(\theta_jy-\theta_j^2/2)}
 {\sum_k w_k\exp(\theta_ky-\theta_k^2/2)}.
\]
The omitted common factor $\phiN(y)$ cancels.  Define the posterior mean,
variance, and cumulants
\[
 \mu(y)=\sum_jp_j(y)\theta_j,\quad
 V(y)=\sum_jp_j(y)\{\theta_j-\mu(y)\}^2,\quad
 \kappa_s(y)=\text{$s$th posterior cumulant}.
\]

\begin{proposition}[Posterior flow and risk curvature]
\label{prop:flow}
For a finite bounded prior, all functions above are analytic and
\[
 \mu'=V,\qquad V'=\kappa_3,\qquad \kappa_s'=\kappa_{s+1}.
\]
With $q(t,z)=\mu(t+z)-t$ and $Z\sim\Normal(0,1)$,
\begin{align}
 \risk(t,\delta_\pi)&=\E q(t,Z)^2,\\
 \risk'(t,\delta_\pi)&=2\E\{q(t,Z)(V(t+Z)-1)\},\\
 \risk''(t,\delta_\pi)&=
 2\E\{(V(t+Z)-1)^2+q(t,Z)\kappa_3(t+Z)\}.
 \label{eq:curvature}
\end{align}
\end{proposition}

\begin{proof}
The probabilities form a finite natural exponential family with natural
parameter $y$.  Successive derivatives of
\[
 \log\!\left(\sum_jw_j e^{\theta_jy-\theta_j^2/2}\right)
\]
produce the cumulants.  Further,
$\partial_tq=V(t+z)-1$ and
$\partial_tV=\kappa_3$.  Two differentiations of $\E q^2$ give the displayed
formulas.  Bounded support supplies an integrable uniform bound for each
derivative.
\end{proof}

The only sign-indefinite product in \eqref{eq:curvature} is $q\kappa_3$.
The following representation couples its factors.

\begin{lemma}[Common-mean representation and enclosure]
\label{lem:common-mean-bound}
Let the posterior be supported on $[\alpha,\beta]$.  If $V>0$, define
\[
 \nu=\frac{\E\{\Theta(\Theta-\mu)^2\mid y\}}{V}.
\]
Then $\nu\in[\alpha,\beta]$ and
\[
 \kappa_3=V(\nu-\mu).
\]
If $t\in[a,b]$, $\mu\in M=[\mu_L,\mu_U]$, and $0\le V\le V_U$, then
\begin{align}
 q\kappa_3
 &\le V_U\max_{\mu\in M}
 \bigl\{(\mu-a)_+(\beta-\mu),\,
        (b-\mu)_+(\mu-\alpha)\bigr\}, \label{eq:joint-upper}\\
 q\kappa_3
 &\ge -V_U\max_{\mu\in M}
 \bigl\{(\mu-a)_+(\mu-\alpha),\,
        (b-\mu)_+(\beta-\mu)\bigr\}. \label{eq:joint-lower}
\end{align}
Each maximum is either an endpoint value or the value at the vertex of one
quadratic.
\end{lemma}

\begin{proof}
The weights proportional to
$p_j(\theta_j-\mu)^2$ are nonnegative and sum to $V$, hence $\nu$ is their
mean and lies in the convex hull $[\alpha,\beta]$.  Expanding
$\E\{\Theta(\Theta-\mu)^2\}= \kappa_3+\mu V$ proves the identity.
For a positive product, if $\mu\ge t$ choose the largest feasible
$\nu=\beta$; if $\mu\le t$ choose $\nu=\alpha$.  Maximizing over
$t\in[a,b]$ gives \eqref{eq:joint-upper}.  The opposite choices of $\nu$
give the most negative product and \eqref{eq:joint-lower}.  The case $V=0$
has $\kappa_3=0$ and follows by continuity.
\end{proof}

\paragraph{Equality and degeneracy.}
Equality in the identity holds by construction.  An envelope endpoint can be
attained only when the posterior squared-deviation weights put their mean
$\nu$ at the relevant support endpoint and when $V=V_U$; $V=0$ collapses all
terms to zero.  These conditions are useful diagnostics but are not assumed
by the solver.

\paragraph{Scope of the enclosure.}
The relaxation separates $V$ from $\mu$ through $V_U$ and encloses, rather than
identifies, their exact joint range.  The ablation compares it with a Cartesian product of independent intervals
for $q$ and $\kappa_3$. The results in
Section~\ref{sec:results} measures the remaining over-enclosure.
